\chapter{Introduction}\label{ch:introduction}
The growing interconnection of systems and the increasing use of automation are contributing to a sharp rise in system complexity, both across and within individual systems \citep{incose2035}. At a system level, independently developed systems are increasingly engineered to cooperate to achieve objectives that exceed their individual capabilities. Such arrangements are described as systems of systems (SoS): large, distributed collections of autonomous systems that collaborate toward higher-level goals while retaining independent operation \citep{maier1998architecting}. Because SoS constituents may join, leave, or reorganize in response to changing objectives, their collective behaviour evolves and gives rise to outcomes that individual systems cannot achieve alone \citep{nielsen2015systems}.

Within individual systems, internal complexity is also increasing as modern software systems collect, interpret, and act on streams of observations. Many contemporary systems therefore operate according to event-driven semantics, in which internal state is updated, and higher-level behaviour is triggered by the arrival of discrete events over time \citep{cugola2012processing}. Digital twins (DTs) exemplify this paradigm. A DT is a digital representation of a system that mirrors and may actuate on its physical counterpart through ongoing bidirectional data exchange \citep{kritzinger2018digital}. These exchanges consist of observations and control updates that drive state updates and computation within the digital representation. These event-driven interactions enable situational awareness \citep{pires2019digital, zhang2024digital}, support analysis and forecasting, and provide services such as simulation-based optimization and decision support \citep{leng2021digital}.

The convergence of system-level coordination in SoS with the real-time representations provided by digital twins gives rise to a new class of systems referred to in this thesis as \emph{systems of twinned systems (SoTS)}.

\begin{definitionframe}
A System of Twinned Systems (SoTS) comprises digitally twinned systems, i.e., digital and physical twins, organized by system of systems (SoS) principles, in which digitally twinned systems may act as autonomous constituents and collaborate to achieve complex goals.
\end{definitionframe}

SoTS extend SoS coordination into the digital domain by enabling constituent systems to interact through their digital twins while maintaining operational autonomy over their physical counterparts. Because digital twins maintain and evolve their internal representations through the continuous exchange of data observations, coordination and decision-making in SoTS depend on the exchange and interpretation of event streams across multiple, distributed twins. As a result, SoTS operate fundamentally as event-driven systems at the computational level.

A defining characteristic of SoTS is therefore their reliance on continuous event flows. At the same time, the properties inherited from SoS make uninterrupted data availability difficult to guarantee. Constituents may reconfigure, communication paths may change, sensing quality may fluctuate, and network or platform conditions may vary over time. These dynamics, which are inherent to SoS operation, manifest at the digital level as temporary gaps, delays, or absences in event streams. Digital twin services depend on these streams for maintaining an up-to-date representation of the observed process. Although such conditions are common in realistic deployments, their implications for SoTS behaviour are often insufficiently addressed.

\paragraph{On the reliability of SoTS.}
Reliability, as defined by \citet{avizienis2004basic}, is the ability of a system to continuously deliver correct service. In SoTS, this requirement applies not only to physical constituents but also to DT-based services such as monitoring, simulation, and analytic decision support. These services rely on an up-to-date digital representation of the observed process, as maintained through event streams. When expected observations fail to arrive, this representation becomes incomplete, causing digital behaviour to diverge from physical reality. In interconnected SoTS, such divergence may propagate across multiple twins, degrading the reliability of collective operation.

From a computational perspective, this reliability challenge arises because event-driven reasoning advances only in response to events. When event streams become incomplete, reasoning processes may stall or operate on outdated information, despite the physical system continuing to evolve. Reliability in SoTS is therefore closely tied to how event-driven computation behaves in the presence of missing or delayed data.


\paragraph{On the role of complex event processing.}
Within event-driven systems such as SoTS, complex event processing (CEP) provides the computational mechanisms for interpreting event streams and detecting higher-level situations \citep{luckham2002power}. In SoTS, CEP forms a critical layer through which low-level observations produced by digital twins are transformed into higher-level representations that support coordination and decision-making across systems.

However, conventional CEP engines implicitly assume uninterrupted event arrival. Across automata-based, logic-based, Petri-net, and grammar-based frameworks, reasoning is defined only over events that actually appear in the input stream. While existing uncertainty-aware CEP approaches address noisy values, uncertain timestamps, or rule-level ambiguity, reasoning progresses only when new events arrive and effectively suspends during periods without data \citep{alevizos2017probabilistic}. As a result, current CEP pipelines provide limited support for reasoning about system state when expected observations fail to arrive. A condition that naturally arises in SoTS due to their SoS characteristics.


\paragraph{Thesis focus and approach.}
This thesis addresses reliability limitations in event-driven computation within SoTS that arise from the temporary absence of expected observations. The focus is on maintaining computational reliability under partial data availability.

The proposed approach preserves continuity of event-driven reasoning by explicitly detecting when expected observations fail to arrive and inferring their likely values through state-based prediction. The inferred information is represented as predicted event data and annotated with confidence values that quantify the expected correctness of each estimate, allowing downstream event processing to continue while making uncertainty explicit.

Architecturally, the framework is realized as a preprocessing mechanism deployed within individual digital twins, upstream of complex event processing. Because digital twins form the computational building blocks of SoTS, this design allows reliability-enhancing mechanisms to be applied locally within each twin and composed across SoS configurations. The framework is grounded in the ISO~23247 reference architecture \cite{iso23247_2_2021} for digital twins, which provides a standardized architectural context that is extended to support reliable event-driven computation under temporary data unavailability.


\section*{Research Contributions}

This thesis advances the reliable operation of SoTS through the following contributions:

\textbf{Contribution 1 – Characterizing Reliability Gaps in Systems of Twinned Systems (\chref{ch:literature-review})}  
Provides a systematic analysis of existing research on systems of twinned systems, characterizing their structural, architectural, and behavioural properties. The analysis identifies a recurring gap in existing work: limited support for maintaining reliable operation when expected observations are temporarily unavailable.

\textbf{Contribution 2 – A Formal Event Model and Architecture for Reliable Event-Driven Operation in SoTS (Chapters~\ref{ch:model} and~\ref{ch:architecture})}  
Develops a formal model that characterizes events and the absence of expected observations, together with a system architecture that supports reliable event-driven computation in SoTS under temporary data unavailability. The architecture introduces mechanisms for detecting missing observations, augmenting event streams with inferred events that are annotated with quantified uncertainty, and preserving the continuity of event-driven reasoning.

\textbf{Contribution 3 – Implementation and Evaluation of the Architecture (Chapters~\ref{ch:implementation} and~\ref{ch:results})}  
Implements the proposed architecture within a CEP-based processing pipeline and evaluates its behaviour under simulated data disruptions. The evaluation assesses reconstruction accuracy, confidence behaviour, and the impact on higher-level event detection, demonstrating the feasibility and effectiveness of the proposed architectural approach.

\vspace{1em}
Together, these contributions establish a foundation for reliable reasoning in systems of twinned systems and demonstrate how architectural support for partial data availability can preserve correctness over time. The chapters that follow develop this foundation by first grounding the work in relevant background and existing research (\chref{ch:background}, \chref{ch:literature-review}), then introducing a formal model and architecture for reliable event-driven operation (\chref{ch:model}, \chref{ch:architecture}). The feasibility and impact of the proposed approach are then demonstrated through implementation and evaluation (\chref{ch:implementation}, \chref{ch:results}), followed by a reflection on limitations, implications, and directions for future research (\chref{ch:discussion}, \chref{ch:conclusion}).

